\documentclass[border=5mm]{standalone}
% ==============================================================================
% SETUP.TEX - CONFIGURACIÓN GLOBAL DEL PROYECTO
% ==============================================================================
% Este archivo centraliza todos los paquetes y estilos.
% Debe incluirse en el preámbulo del libro principal y de los archivos standalone.
% \PassOptionsToPackage{gray}{xcolor}
% --- 1. CODIFICACIÓN E IDIOMA ---
% fix-cm habilita las fuentes Computer Modern en tamaños arbitrarios (por debajo
% de 5 pt, entre otros). Debe cargarse antes que fontenc.
\usepackage{fix-cm}
\usepackage[utf8]{inputenc}
\usepackage[T1]{fontenc}
\usepackage[spanish, es-tabla]{babel}  
\usepackage{csquotes} % Recomendado para babel y citas

\usepackage{import}
% mode=image: Usa el PDF pre-compilado, nunca intenta recompilar.
% Debes compilar cada figura manualmente antes de compilar el libro.
\usepackage[mode=image, subpreambles=false]{standalone}

% --- 2. MATEMÁTICAS Y CIENCIAS ---
\usepackage{amsmath, amssymb, amsfonts, mathtools}
\usepackage{enumitem}

% LaTeX trae \DeclareMathSizes{5}{5}{5}{5}, así que a 5 pt (\tiny) los subíndices
% salen del mismo tamaño que la base: en las etiquetas de figuras el M de
% $\omega_M$ queda desproporcionado. Se restablece la proporción habitual.
\DeclareMathSizes{5}{5}{3.7}{3}

% --- 3. DISEÑO Y GEOMETRÍA --- 
% Unificamos los márgenes para todo el documento
\usepackage[left=2.5cm,right=2.5cm,top=3cm,bottom=3cm]{geometry}
\makeatletter
\ifstandalone % Evita que geometry dañe el recorte de las figuras bajándolas 3cm
  \@ifclasswith{standalone}{crop=false}{
    % Es un capítulo/sección: Dejamos que geometry actúe.
  }{
    % Es una figura: Neutralizamos geometry para que no las recorte mal.
    \geometry{pass}
  }
\fi
\makeatother
\usepackage{parskip} % Espaciado entre párrafos sin sangría
\usepackage{float}   % Para usar [H] en figura

% Ajustes de flotantes para evitar espacios verticales exagerados
\renewcommand{\topfraction}{0.95}      % Permite que las figuras ocupen hasta el 95% superior
\renewcommand{\bottomfraction}{0.95}   % Permite que las figuras ocupen hasta el 95% inferior
\renewcommand{\textfraction}{0.05}     % Permite hojas con tan solo un 5% de texto
\renewcommand{\floatpagefraction}{0.8} % Requiere que una página de solo figuras esté 80% llena
\setcounter{topnumber}{4}
\setcounter{bottomnumber}{4}
\setcounter{totalnumber}{10}

% --- 4. GRÁFICOS Y TABLAS ---
\usepackage{graphicx}
\usepackage{booktabs} % Tablas profesionales
\usepackage{caption}  % Control de captions y \captionof
\usepackage{tikz}
\usetikzlibrary{arrows.meta, calc, angles, quotes, babel, backgrounds}
\usepackage{pgfplots}
\usepgfplotslibrary{groupplots}
\usepgfplotslibrary{external}
\usepgfplotslibrary{patchplots}
\pgfplotsset{compat=1.18}
\usepackage[american, straightvoltages]{circuitikz}

% --- 5. COLORES Y CAJAS ---

\usepackage[table]{xcolor}
\usepackage{tcolorbox}

% Definición de colores corporativos/estilo
\definecolor{utnblue}{RGB}{0, 51, 102}
\definecolor{codegreen}{rgb}{0,0.6,0}
\definecolor{codegray}{rgb}{0.5,0.5,0.5}
\definecolor{codepurple}{rgb}{0.58,0,0.82}
\definecolor{backcolour}{rgb}{0.96,0.96,0.96}

% --- 6. ENTORNOS PERSONALIZADOS (CAJAS) ---

% Caja "Resultado" (Azul Corporativo) — Definiciones, fórmulas clave, resultados importantes.
% Uso: \begin{resultado}[Fórmula de Euler] ... \end{resultado}
\newtcolorbox{resultado}[1][]{
    colback=utnblue!5!white,
    colframe=utnblue,
    title=\textbf{#1},
    fonttitle=\bfseries,
    sharp corners,
    boxrule=0.5mm
}

% Caja "Nota" (Gris) — Advertencias, aplicaciones, contexto secundario.
% Uso: \begin{nota}[Cuidado con la calculadora] ... \end{nota}
% Uso: \begin{nota} ... \end{nota}  → título por defecto "Nota"
\newtcolorbox{nota}[1][Nota]{
    colback=codegray!5!white,
    colframe=codegray,
    title=\textbf{#1},
    fonttitle=\bfseries,
    sharp corners,
    boxrule=0.5mm
}

% Caja inline para resaltar un valor y enlazarlo con su otra aparición en una tabla
% o en una cuenta: mismo color, mismo valor.
% Uso: \marcavalor{$4$}  o  \marcavalor[codegreen]{$4$}
\newtcbox{\marcavalor}[1][utnblue]{
    on line,
    boxsep=1pt, left=2pt, right=2pt, top=1pt, bottom=1pt,
    colback=#1!10!white,
    colframe=#1,
    boxrule=0.4pt,
    arc=2pt
}

% --- 7. CÓDIGO FUENTE (LISTINGS) ---
\usepackage{listings}

\lstdefinestyle{miestilo}{
    backgroundcolor=\color{backcolour},   
    commentstyle=\color{codegreen},
    keywordstyle=\color{magenta},
    numberstyle=\tiny\color{codegray},
    stringstyle=\color{codepurple},
    basicstyle=\ttfamily\footnotesize,
    breakatwhitespace=false,         
    breaklines=true,                 
    captionpos=b,                    
    keepspaces=true,                 
    numbers=left,                    
    numbersep=5pt,                  
    showspaces=false,                
    showstringspaces=false,
    showtabs=false,                  
    tabsize=2,
    frame=single,                    
    rulecolor=\color{codegray}
}

\lstset{style=miestilo}
% Soporte para tildes y ñ en bloques de código
\lstset{literate={ñ}{{\~n}}1 {á}{{\'a}}1 {é}{{\'e}}1 {í}{{\'i}}1 {ó}{{\'o}}1 {ú}{{\'u}}1}

% Operadores matemáticos personalizados

% Vector en sentido discreto (lista finita de coordenadas): negrita + flecha.
% Uso: \vect{v}, \vect{e}_k, \vect{0}.
\newcommand{\vect}[1]{\vec{\mathbf{#1}}}

% Fecha de versión y comando para capítulos standalone
\providecommand{\verdate}{\the\year.\ifnum\month<10 0\fi\the\month.\ifnum\day<10 0\fi\the\day}
\newcommand{\chapterversion}{%
    \vspace{-1.5cm}\begin{flushright}\small\textit{\color{codegray}Versión~\verdate}\end{flushright}\vspace{0.5cm}%
}

% --- 8. HIPERVÍNCULOS (DEBE SER EL ÚLTIMO PAQUETE) ---
\usepackage[colorlinks=true, linkcolor=utnblue, urlcolor=utnblue, citecolor=utnblue]{hyperref}

% --- 9. REFERENCIAS CRUZADAS ENTRE CAPÍTULOS STANDALONE ---
% Cuando se compila un capítulo standalone, xr-hyper lee los .aux de los
% otros capítulos (si existen) para resolver \ref{} a labels externos.
% En la compilación del libro completo este bloque no se activa.
\ifstandalone
  \usepackage{xr-hyper}
  \makeatletter
  \newcommand{\xrchap}[1]{%
    \edef\xrc@self{\detokenize{Capitulo#1}}%
    \edef\xrc@job{\jobname}%
    \ifx\xrc@self\xrc@job\else
      \IfFileExists{../#1capitulo/Capitulo#1.aux}%
        {\externaldocument{../#1capitulo/Capitulo#1}}{}%
    \fi
  }
  \makeatother
  \xrchap{01}\xrchap{02}\xrchap{03}\xrchap{04}
  \xrchap{05}\xrchap{06}\xrchap{07}\xrchap{08}
  \xrchap{09}\xrchap{10}\xrchap{11}\xrchap{12}
  \xrchap{13}
\fi
% \PassOptionsToPackage{gray}{xcolor}

\begin{document}
    \begin{tikzpicture}[>=Latex, scale=1.2, font=\small]
        % Estilos
        \tikzset{
            axis_style/.style={->, thick, black},
            const_r_style/.style={utnblue, thick},
            const_x_style/.style={red!70!black, thick}
        }
        
        % --- PANEL A: PLANO DE IMPEDANCIAS ---
        \begin{scope}[shift={(-4,0)}]
            \node[anchor=south] at (0, 2.8) {\textbf{Plano de Impedancias ($z$)}};
            
            % Fondo: Semiplano derecho (R > 0)
            \fill[cyan!15] (0, -2.5) rectangle (2.5, 2.5);
            
            % Ejes
            \draw[axis_style] (-0.5,0) -- (2.8,0) node[right] {$R$ (Resistencia)};
            \draw[axis_style] (0,-2.5) -- (0,2.5) node[above] {$X$ (Reactancia)};
            
            % Líneas de R constante (Verticales - Azules)
            \foreach \r in {0.5, 1, 1.5, 2} {
                \draw[const_r_style, opacity=0.6] (\r, -2.5) -- (\r, 2.5);
            }
            
            % Líneas de X constante (Horizontales - Rojas)
            \foreach \x in {-2, -1, 0, 1, 2} {
                \draw[const_x_style, opacity=0.6] (0, \x) -- (2.5, \x);
            }
            
            % Etiquetas
            \node[utnblue, anchor=south] at (1, -2.8) {$R = \text{const.}$};
            \node[red!70!black, anchor=south] at (2.3, -2.8) {$X = \text{const.}$};
        \end{scope}

        % Flecha de Transformación
        \draw[->, ultra thick, gray!60] (-0.8, 0.8) to[bend left=15] node[midway, above, black] {$\Gamma = \frac{z-1}{z+1}$} (1.5, 0.8);

        % --- PANEL B: CARTA DE SMITH ---
        \begin{scope}[shift={(4.5,0)}]
            \node[anchor=south] at (0, 2.8) {\textbf{Carta de Smith ($\Gamma$)}};
            
            % Fondo: Disco unitario
            \fill[cyan!15] (0,0) circle (2);
            
            % Ejes
            \draw[axis_style] (-2.5,0) -- (2.5,0) node[right] {Re$(\Gamma)$};
            \draw[axis_style] (0,-2.5) -- (0,2.5) node[above] {Im$(\Gamma)$};
            
            % Círculo unitario (Borde)
            \draw[utnblue!80!black, thick] (0,0) circle (2);
            
            % --- Círculos de R constante (Azules) ---
            % Centro en (r/(r+1), 0), Radio 1/(r+1)
            % R=0: Centro (0,0), Radio 1 (el borde)
            % R=0.5: Centro (1/3, 0) = (0.33, 0), Radio 2/3 = 0.67. Escala x2: (0.67, 0), R=1.33
            % R=1: Centro (0.5, 0), Radio 0.5. Escala x2: (1, 0), R=1
            % R=2: Centro (2/3, 0), Radio 1/3. Escala x2: (1.33, 0), R=0.67
            \begin{scope}
                \clip (0,0) circle (2);
                \draw[const_r_style, opacity=0.7] (0.67, 0) circle (1.33);   % R=0.5
                \draw[const_r_style, opacity=0.7] (1, 0) circle (1);       % R=1
                \draw[const_r_style, opacity=0.7] (1.33, 0) circle (0.67); % R=2
            \end{scope}
            
            % --- Arcos de X constante (Rojos) ---
            % Centro en (1, 1/x), Radio 1/|x|
            % X=1: Centro (1, 1), Radio 1. Escala x2: (2, 2), R=2
            % X=0.5: Centro (1, 2), Radio 2. Escala x2: (2, 4), R=4
            % X=-1: Centro (1, -1), Radio 1. Escala x2: (2, -2), R=2
            \begin{scope}
                \clip (0,0) circle (2);
                \draw[const_x_style, opacity=0.7] (2, 2) circle (2);    % X=+1
                \draw[const_x_style, opacity=0.7] (2, 4) circle (4);    % X=+0.5
                \draw[const_x_style, opacity=0.7] (2, -2) circle (2);   % X=-1
                \draw[const_x_style, opacity=0.7] (2, -4) circle (4);   % X=-0.5
            \end{scope}
            
            % Etiquetas
            \node[utnblue, anchor=north west, scale=0.8] at (0.8, 1.0) {$R=1$};
            \node[red!70!black, anchor=south, scale=0.8] at (0.5, 1.7) {$X=1$};
            
            % Punto de ejemplo (Matched: z=1 -> Gamma=0)
            \fill[codepurple] (0,0) circle (2pt) node[anchor=north east] {\footnotesize Matched};
            
            % Punto Open Circuit (z=inf -> Gamma=1)
            \fill[black] (2,0) circle (2pt) node[anchor=north west] {\footnotesize Open};
            
            % Punto Short Circuit (z=0 -> Gamma=-1)
            \fill[black] (-2,0) circle (2pt) node[anchor=north east] {\footnotesize Short};
        \end{scope}

    \end{tikzpicture}
\end{document}
